\chapter{Corporate Finance Core}\label{ch:corporate-finance-core}

% Scope: time value of money, NPV, the discount rate as hurdle, the growing perpetuity.
% This chapter demonstrates the per-concept template applied to the single most
% reused formula in the book (its marketing twin lives in \cref{ch:customer-economics}).

\section{The Growing Perpetuity}\label{sec:growing-perpetuity}

\begin{decision}
What is a stream of cash worth today if it grows at a constant rate \(g\) forever?
This one model underlies DCF terminal value and customer lifetime value alike.
\end{decision}

\begin{intuition}
A dollar next year is worth less than a dollar today --- today's dollar can be invested.
Growth partly offsets that discount. When the discount rate beats the growth rate,
an infinite stream still sums to a finite number.
\end{intuition}

\begin{derivation}
Sum the discounted cash flows, recognising a geometric series with ratio
\(\frac{1+g}{1+r}<1\):
\begin{equation}\label{eq:gordon}
  PV \;=\; \sum_{t=1}^{\infty}\frac{D\,(1+g)^{t-1}}{(1+r)^{t}}
      \;=\; \frac{D}{\,r-g\,}.
\end{equation}
The denominator \((r-g)\) is the net drag: discount pressure minus the growth lift.
\end{derivation}

\begin{breaks}
The series converges \emph{only} when \(r>g\). If \(g\ge r\) the denominator turns
zero or negative and \cref{eq:gordon} returns nonsense. This is the single most
common valuation error --- assuming a long-run growth rate at or above the discount rate.
\end{breaks}

\begin{worked}
A subscriber cohort throws off \money{1250000} next year at a \qty{42}{\percent}
contribution margin, growing \qty{3}{\percent} per year, discounted at \qty{11}{\percent}:
\[
  PV \;=\; \frac{\money{1250000}\times 0.42}{0.11-0.03}\;\approx\;\money{6562500}.
\]
\end{worked}

\begin{trap}
Reusing the same high \(g\) for ten explicit years \emph{and} the terminal value
double-counts growth. The terminal \(g\) must reflect only the steady state ---
usually at or below long-run GDP.
\end{trap}

\begin{deeper}
Relax constant growth and the model becomes a two-stage or H-model; relax constant
churn in its marketing twin and customer value needs a shifted-beta-geometric survival
model. The clean formula is the \emph{floor} of the analysis, not the ceiling.
\end{deeper}

The identical algebra reappears as customer lifetime value in
\cref{ch:customer-economics}: rename \(D\) the annual customer margin and \(g\) its
growth rate. \textcite{brealey2020} develop the corporate-finance side in full.
